\documentclass[12pt]{article}
\usepackage[margin=1in]{geometry}
\title{Quarterly Performance Report}
\author{Analytics Department}
\date{Q1 2026}
\begin{document}
\maketitle
\tableofcontents
\newpage

\section{Executive Summary}
This report summarizes the key performance indicators for Q1 2026. Revenue
increased by 15\% year-over-year, while operating costs decreased by 3\%.

\section{Financial Overview}
\subsection{Revenue Analysis}
Total revenue reached \$4.2 million, driven primarily by enterprise subscriptions.
The SaaS segment contributed 68\% of total revenue, up from 61\% in Q4 2025.

\subsection{Cost Structure}
Operating expenses totaled \$2.8 million. The largest categories were:
personnel (52\%), infrastructure (23\%), and marketing (15\%).

\subsection{Profitability}
Net income was \$1.1 million, representing a 26\% margin. This is a significant
improvement from the 19\% margin achieved in Q1 2025.

\section{Product Metrics}
\subsection{User Growth}
Monthly active users grew to 142,000, a 28\% increase from the previous quarter.
Daily active users averaged 67,000.

\subsection{Engagement}
Average session duration increased from 8.2 minutes to 11.4 minutes. Feature
adoption rates for the new dashboard exceeded 73\%.

\section{Outlook}
Based on current trends, we project Q2 revenue of \$4.8 million and continued
margin expansion. Key risks include rising infrastructure costs and competitive
pressure in the mid-market segment.
\end{document}